\scanpage{154}
\mspage{28}
% archive sheet number: 153

\uline{Dém} Tout a été prouvé\unsure{,|\}} si on se rappelle que les
objets de $M^r$ sont $r^{\pi_0}$-accessibles (cf. 3.2).

\deleted{\uline{Corollaire} \unsure{La conclusion|Les conclusions} de stabilité \unsure{reste valable|restent valables},
lorsque $s$ est}

4. Compléments sur les structures de Quillen.

\uline{Définition 4.1} Soit $M$ une catégorie. On appelle
\uline{précouple de Quillen} dans $M$ un couple
$(\Phi,\Psi)$ de parties de $\mathrm{Fl}\,M$, satisfaisant les
deux conditions suivantes

a) On a $\Phi\longleftrightarrow\Psi$, i.e. $\Psi\subset\Phi_*$, ou ce qui revient
au $=$, $\Phi\subset\Psi^*$ \deleted{\unsure{et}} -- soit encore : Pour
tout carré commutatif dans $M$

\diagram{\begin{tikzcd}[ampersand replacement=\&,column sep=large]
\& A \arrow[r,"u"] \arrow[d,dash,"i"',"{\unsure{\text{---}|\downarrow}}"] \& X \arrow[d,"p"]\\
\text{(4.1.1.)} \& B \arrow[r,dash,"v"'] \arrow[ur,dotted,"w"] \& Y
\end{tikzcd}}
avec $i\in\Phi$, $p\in\Psi$, il existe $w$ telle que
$wi=u,\ pw=v$.
\margin{\unsure{i.e.|.2.}}

b) Toute $f\in\mathrm{Fl}\,M$ se factorise en
\[
f=pi,\quad\text{avec }i\in\Phi,\ p\in\Psi.
\]

On écrit
\[
\Phi\mathrel{\overset{\vphantom{I}}{\longleftrightarrow}\mkern-19mu\mid\mkern9mu}\Psi
\tag*{(4.1.2)}
\]
pour indiquer que $(\Phi,\Psi)$ est un \inserted{pré-}couple
de Quillen.

On \deleted{appelle} dit que le couple de Quillen $(\Phi,\Psi)$
est \uline{clos}, si $\Phi$ est clos à gauche, $\Psi$ clos à droite.
On va voir plus bas \inserted{(voir \unsure{page|pag.} 31)} que cela implique (donc
équivaut à)
\[
\Phi=\Psi^*,\qquad\Psi=\Phi_*.
\tag*{(4.1.3.)}
\]
